\documentclass[runningheads]{llncs}

\usepackage[T1]{fontenc}
\usepackage[utf8]{inputenc}
\usepackage{graphicx}
\usepackage[color]{coqdoc}
\usepackage{mathpartir}
\usepackage{amssymb}
\usepackage{newunicodechar}
\usepackage{tcolorbox}
\usepackage{xcolor}
\usepackage[all,color]{xy}
\usepackage{url}
\providecommand{\doi}[1]{\url{https://doi.org/#1}}

\newunicodechar{∀}{\ensuremath{\forall}}
\newunicodechar{∈}{\ensuremath{\in}}
\newunicodechar{∉}{\ensuremath{\notin}}
\newunicodechar{❬}{\ensuremath{\langle}}
\newunicodechar{❭}{\ensuremath{\rangle}}
\newunicodechar{α}{\ensuremath{\alpha}}
\newunicodechar{β}{\ensuremath{\beta}}
\newunicodechar{…}{\ldots}
\newunicodechar{↠}{\ensuremath{\tto}}

\newcommand{\vswap}[3]{(\!\!(#1\ #2)\!\!) #3}
\newcommand{\swap}[3]{(#1\ #2) #3}
\newcommand{\metasub}[3]{\{ #2 := #3 \}#1}
\newcommand{\graft}[3]{\{ #2 /\!\!/ #3 \}#1}
\newcommand{\renv}[3]{\{ #2 \leadsto #3 \}#1}
\newcommand{\esub}[3]{[ #2 := #3 ] #1}
\newcommand{\tto}{\relbar\joinrel\twoheadrightarrow}

\newcommand{\flavio}[1]{{\color{red}#1}}

\begin{document}  
\title{A Nominal Formalization of the Confluence of a Calculus with Explicit Substitutions via the Z Property}
\titlerunning{A Nominal Formalization of Confluence via the Z Property}
\author{Flávio L. C. de Moura\inst{1}\orcidID{0000-0002-9390-5751}}

\authorrunning{F. L. C. de Moura}

\institute{Departamento de Ciência da Computação \\
  Universidade de Brasília,
  Brasília, Brazil \\
  \email{flaviomoura@unb.br}}

\maketitle

\begin{abstract}
The Z property is a compact proof technique for confluence that has been used to give elegant confluence proofs for the $\lambda$-calculus and several of its extensions. Previous formalizations of the (compositional) Z property, as well as of the Substitution Lemma for calculi with explicit substitutions, were developed in Rocq. In this work, we recast this framework in a nominal setting, where variables are represented by atoms and binders are handled through name-swapping rather than de Bruijn indices, and formally verify it in the Rocq prover (version 9.1.1). As an application of this nominal framework, we formalize the $\lambda_x$-calculus, an extension of the $\lambda$-calculus with explicit substitutions, and give a fully mechanized proof of its confluence via the compositional Z property. The whole development, which is freely available (\url{https://github.com/flaviodemoura/Zprop-nominal.git}), comprises around $4\,500$ lines of Rocq code and $150$ lemmas and theorems, and provides a reusable nominal toolkit for confluence proofs of calculi with explicit substitutions.
\end{abstract}

\keywords{Nominal syntax, Confluence, Z property, Explicit substitutions, Rocq Prover, Formal verification.}

\input{infra_nom.v.tex}
\input{aeq_equiv_es.v.tex}
\input{lx_nom.v.tex}

\section{Related Work and Conclusion}

\noindent{\bf Related Work.} The Z property was introduced by van Oostrom \cite{oostromDraftYourMind2007} and Dehornoy and van Oostrom \cite{dehornoy2008z} as a simple, self-contained alternative to the classical Tait-Martin-Löf technique for proving confluence\cite{hondaConfluenceProofsLambdaMuCalculi2021,nakazawaCallbyvalue2017,nakazawaPropertyShufflingCalculus2023}, and it has since been used as a benchmark for confluence tools and formalizations \cite{felgenhauerProperty2016,vanoostrom:LIPIcs.FSCD.2021.24}. Nakazawa and Fujita \cite{nakazawaCompositionalConfluenceProofs2016} proposed the compositional Z property, an extension tailored to calculi whose reduction relation splits naturally into an initiating step and a substitution calculus that completes it, and used it to give a pencil-and-paper confluence proof for a calculus with explicit substitutions close to the $\lambda_x$-calculus studied here. The present paper can be seen as a full mechanization, and a nominal reformulation, of that proof strategy.

Our development directly builds on two previous formalizations by (a subset of) the same authors: \cite{fmm2021}, which formalizes the Z and compositional Z properties, and \cite{limaFormalizedExtensionSubstitution2023}, which extends the classical substitution lemma to calculi with explicit substitutions. Both were carried out using a locally nameless representation of binders, built on top of the \texttt{Metalib} library \cite{weirichCoqMetalibPackage}. Here we keep the same underlying atom infrastructure but switch to a nominal representation, in which $\alpha$-equivalence is defined explicitly via name-swapping \cite{gabbayNewApproachAbstract2002,pittsNominalLogicFirst2003} rather than being identified with syntactic equality of de Bruijn-style terms. This choice trades some of the automation that locally nameless representations provide for a syntax that is closer to the informal, textbook presentation of binders, which we found particularly convenient when stating and proving properties of the metasubstitution and of the $\lambda_x$-calculus reduction rules. Related nominal infrastructures have been developed in other proof assistants, notably Nominal Isabelle \cite{urbanGeneralBindingsAlphaEquivalence2012} and a nominal library for Agda \cite{gabbayNominalTechniquesAgda}; unlike these, our formalization does not rely on a dedicated nominal package or on quotient types, but implements swapping and $\alpha$-equivalence directly as inductive definitions over raw terms, in the same spirit as other machine-checked accounts of the Barendregt variable convention \cite{copelloMachinecheckedProofChurchRosser2018,urbanBarendregtsVariableConvention2007,aydemirEngineeringFormalMetatheory2008}.

\smallskip
\noindent{\bf Conclusion.} We presented a nominal formalization, in the Rocq proof assistant (version 9.1.1), of a framework for proving confluence via the (compositional) Z property, together with the supporting theory of $\alpha$-equivalence, swapping, and capture-avoiding metasubstitution for nominal terms with an explicit substitution operator. As an application of this framework, we formalized the $\lambda_x$-calculus and proved that its reduction relation $\to_{lx}$ is confluent, by exhibiting the complete permutation function $P$ and the complete development function $B$ as witnesses of the compositional Z property. To the best of our knowledge, this is the first fully mechanized, nominal proof of confluence for a calculus with explicit substitutions via the (compositional) Z property. The whole development is freely available (\url{https://github.com/flaviodemoura/Zprop-nominal.git}) and comprises around $4\,500$ lines of Rocq code, organized in $150$ lemmas and theorems, which can serve as a reusable toolkit for confluence proofs of other calculi with explicit substitutions.

As future work, we intend to extend the framework with metavariables, which play an important role in the study of unification, and to apply it to other calculi with explicit substitutions from the long list surveyed in the introduction \cite{abadiExplicitSubstitutions1991,roseExplicitSubstitutionNames2011,benaissaLnCalculusExplicit1996,curienConfluencePropertiesWeak1996,munozConfluencePreservationStrong1996,kamareddineExtendingLcalculusExplicit1997a,blooExplicitSubstitutionEdge1999,davidLambdacalculusExplicitWeakening2001,kesnerTheoryExplicitSubstitutions2009a}, including calculi that support composition of substitutions, a feature the $\lambda_x$-calculus deliberately avoids. We also plan a systematic, quantitative comparison between the nominal formalization presented here and the locally nameless formalizations \cite{fmm2021,limaFormalizedExtensionSubstitution2023} it builds upon, in terms of proof size, automation effort, and robustness to changes in the calculus under study.

\bibliographystyle{splncs04}
\bibliography{report} 
\end{document}
